% GENERATED FILE. Edit canonical lessons, then run npm run generate:books.
% canonical-source-hash: fnv1a64:40f27f3e194130d5
% canonical-lessons: GE-C328-sieb, GE-C328-abgiessen, GE-C328-abtropfen, GE-C328-abschmecken, GE-C328-wuerzen

\chapter{Abgießen und abschmecken}
\label{ch:ge-abgie-en-und-abschmecken}
\hlchaptermodality{328}

\begin{chapteropening}
\textbf{By the end of this chapter:} \emph{I can describe finishing a dish: a sieve, to drain, to drip dry, to taste for seasoning and to season.}

Complete the last lesson of chapter 328: I can describe finishing a dish: a sieve, to drain, to drip dry, to taste for seasoning and to season.
\end{chapteropening}

\section[das Sieb]{das Sieb — the sieve, the colander}

\label{lesson:GE-C328-sieb}



\begin{quote}
Before the new one: say the German for to brown quickly, then the German for the lid.
\end{quote}

\subsection*{You'll want to know: das Sieb}
\textbf{das Sieb} — \textquotedblleft{}the sieve, the colander\textquotedblright{}.

\begin{quote}
Er goss die Nudeln in ein Sieb.
\end{quote}

\begin{quote}
\emph{He poured the pasta into a colander.}
\end{quote}

Plural: \emph{die Siebe}.

\begin{cousinweb}
\emph{Sieb} is inherited from Old High German \emph{sib}, from Proto-Germanic \emph{sibi} \textquotedblleft{}sieve\textquotedblright{} (EtymDB), the same source as English \emph{sieve}.
\end{cousinweb}

\begin{grammarlens}[title={What you've built}]
One more word for telling, explaining and describing at B1.
\end{grammarlens}

\subsection*{Your turn}
\begin{itemize}
\raggedright
  \item \emph{Say it:} \emph{das Sieb}
  \item \emph{Say it:} \emph{das Sieb}, once more
  \item \emph{Say it:} a sentence of your own with \emph{das Sieb}
  \item \emph{Recall:} say the German for to brown quickly, then the German for the lid, then say \emph{das Sieb} again
\end{itemize}

\begin{tcolorbox}[breakable,colback=teal!4,colframe=teal!35!black,title={Before you move on}]
What does das Sieb mean? (\textquotedblleft{}The sieve, the colander\textquotedblright{}.) Say it once more.
\end{tcolorbox}
\section[abgießen]{abgießen — to drain (by pouring off)}

\label{lesson:GE-C328-abgiessen}



\begin{quote}
Before the new one: say the German for the lid, then the German for the sieve.
\end{quote}

\subsection*{You'll want to know: abgießen}
\textbf{abgießen} — \textquotedblleft{}to drain (by pouring off)\textquotedblright{}.

\begin{quote}
Sie goss die Kartoffeln ab.
\end{quote}

\begin{quote}
\emph{She drained the potatoes.}
\end{quote}

Separable, strong: \emph{goss ab}, \emph{hat abgegossen}.

\begin{cousinweb}
\emph{abgießen} is built on \emph{gießen} \textquotedblleft{}to pour\textquotedblright{}, Old High German \emph{giozan}, Proto-Germanic \emph{geutaną} (EtymDB), with \emph{ab} \textquotedblleft{}off\textquotedblright{}.
\end{cousinweb}

\begin{grammarlens}[title={What you've built}]
One more word for telling, explaining and describing at B1.
\end{grammarlens}

\subsection*{Your turn}
\begin{itemize}
\raggedright
  \item \emph{Say it:} \emph{abgießen}
  \item \emph{Say it:} \emph{abgießen}, once more
  \item \emph{Say it:} a sentence of your own with \emph{abgießen}
  \item \emph{Recall:} say the German for the lid, then the German for the sieve, then say \emph{abgießen} again
\end{itemize}

\begin{tcolorbox}[breakable,colback=teal!4,colframe=teal!35!black,title={Before you move on}]
What does abgießen mean? (\textquotedblleft{}To drain (by pouring off)\textquotedblright{}.) Say it once more.
\end{tcolorbox}
\section[abtropfen]{abtropfen — to drip dry, to drain off}

\label{lesson:GE-C328-abtropfen}



\begin{quote}
Before the new one: say the German for the sieve, then the German for to drain.
\end{quote}

\subsection*{You'll want to know: abtropfen}
\textbf{abtropfen} — \textquotedblleft{}to drip dry, to drain off\textquotedblright{}.

\begin{quote}
Er ließ den Salat gut abtropfen.
\end{quote}

\begin{quote}
\emph{He let the lettuce drip dry.}
\end{quote}

Separable, often with \emph{lassen}. Perfekt with \emph{sein}.

\begin{cousinweb}
\emph{abtropfen} is built on \emph{tropfen} \textquotedblleft{}to drip\textquotedblright{}, Old High German \emph{tropfōn} (EtymDB), with \emph{ab} \textquotedblleft{}off\textquotedblright{}.
\end{cousinweb}

\begin{grammarlens}[title={What you've built}]
One more word for telling, explaining and describing at B1.
\end{grammarlens}

\subsection*{Your turn}
\begin{itemize}
\raggedright
  \item \emph{Say it:} \emph{abtropfen}
  \item \emph{Say it:} \emph{abtropfen}, once more
  \item \emph{Say it:} a sentence of your own with \emph{abtropfen}
  \item \emph{Recall:} say the German for the sieve, then the German for to drain, then say \emph{abtropfen} again
\end{itemize}

\begin{tcolorbox}[breakable,colback=teal!4,colframe=teal!35!black,title={Before you move on}]
What does abtropfen mean? (\textquotedblleft{}To drip dry, to drain off\textquotedblright{}.) Say it once more.
\end{tcolorbox}
\section[abschmecken]{abschmecken — to taste for seasoning}

\label{lesson:GE-C328-abschmecken}



\begin{quote}
Before the new one: say the German for to drain, then the German for to drip dry.
\end{quote}

\subsection*{You'll want to know: abschmecken}
\textbf{abschmecken} — \textquotedblleft{}to taste for seasoning\textquotedblright{}.

\begin{quote}
Am Ende schmeckte sie die Suppe mit Salz ab.
\end{quote}

\begin{quote}
\emph{At the end she seasoned the soup to taste with salt.}
\end{quote}

Separable: \emph{mit Salz und Pfeffer abschmecken}.

\begin{cousinweb}
\emph{abschmecken} is built on \emph{schmecken} \textquotedblleft{}to taste\textquotedblright{}, from Proto-Germanic \emph{smakkijaną} \textquotedblleft{}to taste\textquotedblright{} (EtymDB), with \emph{ab}.
\end{cousinweb}

\begin{grammarlens}[title={What you've built}]
One more word for telling, explaining and describing at B1.
\end{grammarlens}

\subsection*{Your turn}
\begin{itemize}
\raggedright
  \item \emph{Say it:} \emph{abschmecken}
  \item \emph{Say it:} \emph{abschmecken}, once more
  \item \emph{Say it:} a sentence of your own with \emph{abschmecken}
  \item \emph{Recall:} say the German for to drain, then the German for to drip dry, then say \emph{abschmecken} again
\end{itemize}

\begin{tcolorbox}[breakable,colback=teal!4,colframe=teal!35!black,title={Before you move on}]
What does abschmecken mean? (\textquotedblleft{}To taste for seasoning\textquotedblright{}.) Say it once more.
\end{tcolorbox}
\section[würzen]{würzen — to season, to spice}

\label{lesson:GE-C328-wuerzen}



\begin{quote}
Before the new one: say the German for to drip dry, then the German for to taste for seasoning.
\end{quote}

\subsection*{You'll want to know: würzen}
\textbf{würzen} — \textquotedblleft{}to season, to spice\textquotedblright{}.

\begin{quote}
Er würzte das Fleisch mit Pfeffer.
\end{quote}

\begin{quote}
\emph{He seasoned the meat with pepper.}
\end{quote}

\emph{Das Gewürz} — the spice.

\begin{cousinweb}
\emph{würzen} is built on \emph{Würze} \textquotedblleft{}seasoning\textquotedblright{}, from Middle High German \emph{wurz} \textquotedblleft{}herb, root\textquotedblright{} (EtymDB).
\end{cousinweb}

\begin{grammarlens}[title={What you've built}]
One more word for telling, explaining and describing at B1.
\end{grammarlens}

\subsection*{Your turn}
\begin{itemize}
\raggedright
  \item \emph{Say it:} \emph{würzen}
  \item \emph{Say it:} \emph{würzen}, once more
  \item \emph{Say it:} a sentence of your own with \emph{würzen}
  \item \emph{Recall:} say the German for to drip dry, then the German for to taste for seasoning, then say \emph{würzen} again
\end{itemize}

\begin{tcolorbox}[breakable,colback=teal!4,colframe=teal!35!black,title={Before you move on}]
What does würzen mean? (\textquotedblleft{}To season, to spice\textquotedblright{}.) Say it once more.
\end{tcolorbox}
